\chapter{Prima appendice}

\lipsum[1]

\begin{tabular}{cc}
    
\end{tabular}

\begin{itemize}
    \item Hardware (1*A100) / (4*H100)
    \item vLLM Recipe (quantizzazione)
    \item Usage Parameters (0.8) / (1.0 thinking, 0.8)
\end{itemize}

\lipsum[1]

% ============================================================
% Trucco Avanzato per distinguere Chiavi e Valori in YAML
% ============================================================
\newcommand\YAMLkeystyle{\color{yamlkey}\bfseries}
\newcommand\YAMLvaluestyle{\color{yamlvalue}\mdseries}
\newcommand\YAMLcolonstyle{\color{yamlvalue}\mdseries}

\makeatletter
\newcommand\ProcessThreeDashes{\textcolor{yamlcomment}{---}}

\lstdefinelanguage{yaml}{
    sensitive=true,
    morecomment=[l]{\#},
    morestring=[b]",
    morestring=[b]',
    % Gestione dei due punti per separare chiave e valore
    moredelim=**[il][\YAMLcolonstyle{:}\YAMLvaluestyle]{:},
    % Riconoscimento dei tre trattini di inizio YAML
    literate={---}{{\ProcessThreeDashes}}3,
    % Booleani e keyword principali
    keywords={true, false, null, yes, no, on, off},
    keywordstyle=\color{yamlkeyword}\bfseries
}
\makeatother

% ============================================================
% Palette Colori "Modern Clean" (stile VS Code / GitHub)
% ============================================================
\definecolor{yamlvalue}{HTML}{0550AE}       % Blu per le chiavi (prima dei due punti)
\definecolor{yamlkey}{HTML}{24292F}     % Antracite scuro per i valori standard
\definecolor{yamlkeyword}{HTML}{CF222E}   % Rosso/Amaranto per i booleani
\definecolor{yamlcomment}{HTML}{57606A}   % Grigio per i commenti
\definecolor{yamlnumber}{HTML}{8C95A0}    % Grigio chiaro per i numeri di riga
\definecolor{yamlbackground}{HTML}{F6F8FA} % Sfondo grigio chiarissimo
\definecolor{yamlline}{HTML}{D0D7DE}      % Bordo sottile ed elegante

% ============================================================
% Definizione dello Stile Finale
% ============================================================
\lstdefinestyle{yamlstyle}{
    language=yaml,
    % Di base tutto il testo adotta lo stile della CHIAVE, 
    % poi i due punti lo cambieranno in VALUE
    basicstyle=\small\ttfamily\YAMLkeystyle,
    columns=fullflexible,
    keepspaces=true,
    showstringspaces=false,
    breaklines=true,
    breakatwhitespace=true,
    % Numeri di riga
    numbers=left,
    numberstyle=\tiny\ttfamily\color{yamlnumber},
    stepnumber=1,
    numbersep=12pt,
    % Cornice ed estetica del blocco
    backgroundcolor=\color{yamlbackground},
    frame=single,
    rulecolor=\color{yamlline},
    framerule=0.8pt,
    framesep=8pt,
    xleftmargin=24pt,
    xrightmargin=8pt,
    upquote=true
}

% ============================================================
% YAML listing (Uso)
% ============================================================
\iffalse
\lstinputlisting[
    style=yamlstyle,
    caption={Configurazione YAML di vLLM},
    label={lst:vllm-yaml}
]{appendices/vllm_recipe_q3627.yaml}
\fi

\begin{figure}[t]
    \centering
    % --------------------------------------------------------
    % Prima Minipage (Sinistra) - Larghezza: 48% della pagina
    % --------------------------------------------------------
    \begin{minipage}{0.45\textwidth}
        \lstinputlisting[
            style=yamlstyle,
            basicstyle=\small\ttfamily\YAMLkeystyle, % Font leggermente più piccolo per farlo stare nello spazio
            xleftmargin=12pt, % Margini ridotti per l'affiancamento
            caption={Configurazione A},
            label={lst:vllm-yaml-a}
        ]{appendices/vllm_recipe_q3627.yaml}
    \end{minipage}
    \hspace{0.06\textwidth}% Crea lo spazio elastico orizzontale tra le due minipage
    % --------------------------------------------------------
    % Seconda Minipage (Destra) - Larghezza: 48% della pagina
    % --------------------------------------------------------
    \begin{minipage}{0.48\textwidth}
        \lstinputlisting[
            style=yamlstyle,
            basicstyle=\small\ttfamily\YAMLkeystyle, % Font coerente con il primo
            xleftmargin=12pt,
            caption={Configurazione B},
            label={lst:vllm-yaml-b}
        ]{appendices/vllm_recipe_ds4fp.yaml}
    \end{minipage}
\end{figure}

\chapter{Seconda appendice}

\lipsum[1]

% ============================================================
% Palette Colori "Python Modern"
% ============================================================
\definecolor{pykeyword}{HTML}{CF222E}   % Rosso per keyword (def, class, import, return)
\definecolor{pybuiltin}{HTML}{953097}   % Viola per funzioni built-in (print, len, range)
\definecolor{pystring}{HTML}{0A3069}    % Blu scuro per le stringhe
\definecolor{pycomment}{HTML}{57606A}   % Grigio per i commenti
\definecolor{pynumber}{HTML}{005A9C}    % Blu brillante per i numeri
\definecolor{pybg}{HTML}{F6F8FA}        % Sfondo grigio chiarissimo
\definecolor{pyline}{HTML}{D0D7DE}      % Bordo sottile elegante

% ============================================================
% Definizione dello Stile Python
% ============================================================
\lstdefinestyle{pystyle}{
    language=Python,
    % basicstyle=\small\ttfamily\color{black},
    basicstyle=\scriptsize\ttfamily\color{black},
    keywordstyle=\color{pykeyword}\bfseries,
    stringstyle=\color{pystring},
    commentstyle=\color{pycomment}\itshape,
    % Evidenziazione avanzata per numeri, built-in e decoratori
    emph={print,len,range,str,int,float,dict,list,set,tuple,enumerate,zip},
    emphstyle=\color{pybuiltin},
    moredelim=**[is][\color{pynumber}]{£}{£}, % Per evidenziare i numeri se necessario
    % Formattazione blocco e bordi
    backgroundcolor=\color{pybg},
    frame=single,
    rulecolor=\color{pyline},
    framerule=0.8pt,
    framesep=8pt,
    xleftmargin=24pt,
    xrightmargin=8pt,
    % Gestione spazi e righe
    numbers=left,
    numberstyle=\tiny\ttfamily\color{pycomment},
    numbersep=12pt,
    breaklines=true,
    breakatwhitespace=true,
    showstringspaces=false,
    upquote=true
}

\lstinputlisting[
    style=pystyle,
    caption={BotNet Template},
    label={lst:py}
]{appendices/botnet_template.py}

\chapter{Terza appendice}